\begin{statementbox}
\textbf{\textcolor{CStatement}{条件}}
\begin{description}[style=nextline, leftmargin=2.8em, labelsep=0.5em, itemsep=0.35em]
    \item[\textbf{\textcolor{CStatement}{条件 1（实数定义）：}}] 设 $x$ 为任意实数。
\end{description}

\textbf{\textcolor{CStatement}{结论}}
\begin{description}[style=nextline, leftmargin=2.8em, labelsep=0.5em, itemsep=0.45em]
    \item[\textbf{\textcolor{CStatement}{结论一（切线放缩一）：}}]
    \textbf{\textcolor{CStatement}{直观描述：}} 指数函数在 $x=0$ 处的切线不等式。
    \[
    e^{x} \geq x+1
    \]

    \item[\textbf{\textcolor{CStatement}{结论二（切线放缩二）：}}]
    \textbf{\textcolor{CStatement}{直观描述：}} 指数函数在 $x=1$ 处的切线不等式。
    \[
    e^{x} \geq e x
    \]

    \item[\textbf{\textcolor{CStatement}{结论三（二阶放缩）：}}]
    \textbf{\textcolor{CStatement}{直观描述：}} 当 $x$ 非负时的二阶多项式下界。
    \[
    e^{x} \geq 1+x+\frac{x^{2}}{2} \quad (x \geq 0)
    \]

    \item[\textbf{\textcolor{CStatement}{推广结论（等价变形）：}}]
    \textbf{\textcolor{CStatement}{直观描述：}} 结论二的等价形式，便于变量替换。
    \[
    e^{x-1} \geq x
    \]
\end{description}

\textbf{\textcolor{CStatement}{等号/取等条件：}} 结论一等号成立当且仅当 $x=0$；结论二等号成立当且仅当 $x=1$；结论三等号成立当且仅当 $x=0$；推广结论等号成立当且仅当 $x=1$。
\end{statementbox}
